\scanpage{35}
\uline{N.B.} On dit que, pour un $T$ \uline{local} sur $S$ et $\eta\in F(T)$, $(T,\eta)$ est \uline{modulaire} \inserted{\uline{local}} si

$1^\circ)$ pour $\forall T'$ local sur $S$, l'application
\[
\operatorname{Hom}_{\mathrm{loc},S}(T',T)\longrightarrow F(T')
\]
définie par $\eta$ est une \deleted{application} bijection du premier ens sur le sous-ens du second formé des $\eta'\in F(T')$ tels que $\eta'(t')\in F(k(t'))$ soit dans
\[
\operatorname{Im}\bigl[\operatorname{Hom}_k(k(t'),k(t))\longrightarrow F(k(t'))\bigr],
\]
où la dernière flèche est définie par $\eta(t)\in F(k(t))$.
\margin{$t,t'$ les pts fermés de $T,T'$}

\margin{\figure{\begin{tikzpicture}[x=1cm,y=1cm]
\draw (0,2.8) .. controls (1.7,2.6) and (.8,1.2) .. (2.5,1.2) .. controls (3.5,1.2) and (2.9,-.8) .. (2.6,-1.7) .. controls (2.4,-2) and (1.6,-1.7) .. (1.3,-1.8);
\node at (.45,2.1) {$t\in T$};\node at (.4,.65) {$s\in S$};\node at (2.2,.65) {$T'\ni t'$};
\end{tikzpicture}}}
\margin{\diagram{\begin{tikzcd}[column sep=small,row sep=small]
\deleted{\unsure{L}} & B \deleted{\arrow[l]} \arrow[r] & L \arrow[dr,dashed,leftrightarrow] & \\
& & k \arrow[u] \arrow[r] & L' \\
\deleted{k}\deleted{\arrow[uu]} & A\deleted{\arrow[l]}\arrow[uu]\arrow[ur]\arrow[rr] & & B'\arrow[u]\deleted{\arrow[r]}\deleted{\arrow[d]} & \deleted{\unsure{R'|K'}} \\
& & & \deleted{\unsure{L'|L}} &
\end{tikzcd}}}

$2^\circ)$ Pour $\forall$ diagramme commutatif de corps \textcircled{\unsure{6|a}}, avec $L=k(t)$, $L'=k(t')$, $k\mathbin{-}L$ et $k\mathbin{-}L'$ les homs canoniques, et $\forall\beta'\in F(L')$ tel que $F(j')(\beta')=F(j)(\beta)$ [où $\beta=\eta(t)$], on a
\[
\left\{\begin{array}{l}
j'(L')\supset L\qquad(\text{d'où }\lambda:L\mathbin{-}L')\\
\beta'=F(\lambda)(\beta)
\end{array}\right.
\]
\margin{\textcircled{\unsure{6|a}}\diagram{\begin{tikzcd}
\eta(t)=\beta\quad L \arrow[r,"j"] \arrow[dr,dashed,"\lambda"] & M \\
k\arrow[u,"i"]\arrow[r,"i'"'] & L'\arrow[u,"j'"']\quad\beta'=\unsure{1|i}
\end{tikzcd}}}

[N.B. on \deleted{\unsure{exprime que}} \inserted{exprimera $2^\circ)$ en} disant que $L$ est un \uline{corps de définition} pour $\beta$]
\mspage{34}
